\documentclass[11pt]{article}
\usepackage[a4paper,margin=1in]{geometry}
\usepackage{fontspec}
\usepackage[boldfont=false]{xeCJK}
\setCJKmainfont{FandolSong-Regular.otf}
\usepackage{amsmath}
\usepackage{amssymb}
\usepackage{array}
\usepackage{booktabs}
\usepackage{graphicx}
\usepackage{adjustbox}
\usepackage{enumitem}
\setlist[itemize]{topsep=2pt,partopsep=0pt,parsep=0pt,itemsep=2pt,leftmargin=1.4em}
\setlist[enumerate]{topsep=2pt,partopsep=0pt,parsep=0pt,itemsep=2pt,leftmargin=1.6em}
\usepackage{parskip}
\usepackage{fvextra}
\fvset{breaklines=true,breakanywhere=true,fontsize=\small}
\setlength{\parindent}{0pt}

\begin{document}

\begin{center}
  {\LARGE\bfseries Nitrogen and sulfur}\\[0.35em]
  {\large A Level Chemistry}
\end{center}
\vspace{0.6em}

\section*{Why nitrogen is unreactive}

Nitrogen gas, $\text{N}_2$, makes up most of the air but reacts with very little. There are two reasons:

\begin{itemize}
\item the two nitrogen atoms are joined by a \textbf{triple bond} 三键, which has a very high \textbf{bond energy} 键能. A lot of energy is needed to break it.

\item the molecule has no \textbf{polarity} 极性 — it is perfectly symmetrical, so nothing pulls other molecules towards it.

\end{itemize}

\section*{Ammonia and the ammonium ion}

\subsection*{Basicity of ammonia}

The \textbf{basicity} 碱性 of ammonia (its ability to act as a base) comes from the \textbf{lone pair} 孤对电子 of electrons on the nitrogen atom. Using the Brønsted–Lowry theory, ammonia is a base because this lone pair can accept a \textbf{proton} 质子 ($\text{H}^+$):

\[
\text{NH}_3 + \text{H}^+ \rightarrow \text{NH}_4^+
\]

\subsection*{The ammonium ion}

When the lone pair forms a bond to $\text{H}^+$, it makes the \textbf{ammonium ion} 铵离子, $\text{NH}_4^+$. Because both shared electrons came from the nitrogen, this new bond is a \textbf{coordinate bond} 配位键. The ion has four identical N–H bonds and a tetrahedral shape.

\subsection*{Displacement of ammonia from its salts}

If you warm an ammonium salt with a base (such as sodium hydroxide), you push out ammonia gas. This is an acid–base \textbf{displacement} 置换:

\[
\text{NH}_4\text{Cl} + \text{NaOH} \rightarrow \text{NaCl} + \text{NH}_3 + \text{H}_2\text{O}
\]

The sharp smell of ammonia, and damp red litmus turning blue, is a test for an ammonium salt.

\section*{Oxides of nitrogen and air pollution}

\subsection*{Where they come from}

Oxides of nitrogen ($\text{NO}$ and $\text{NO}_2$, together called $\text{NO}_x$) come from two sources:

\begin{itemize}
\item \textbf{natural}: lightning gives enough energy for nitrogen and oxygen in the air to combine.

\item \textbf{man-made}: the high temperature inside an \textbf{internal combustion engine} 内燃机 makes nitrogen and oxygen react:

\end{itemize}

\[
\text{N}_2 + \text{O}_2 \rightarrow 2\text{NO}
\]

\subsection*{Removing them from car exhaust}

A \textbf{catalytic converter} 催化转化器 cleans the \textbf{exhaust gases} 尾气. It lets the harmful gases react together to form harmless ones:

\[
2\text{NO} + 2\text{CO} \rightarrow \text{N}_2 + 2\text{CO}_2
\]

\subsection*{Photochemical smog}

In sunlight, $\text{NO}$ and $\text{NO}_2$ react with unburned \textbf{hydrocarbons} to form peroxyacetyl nitrate (PAN). PAN is a harmful part of \textbf{photochemical smog} 光化学烟雾, the brown haze seen over busy cities.

\subsection*{Acid rain}

The oxides of nitrogen also help make \textbf{acid rain} 酸雨 in two ways:

\begin{itemize}
\item directly: $\text{NO}_2$ dissolves in rain to form nitric acid.

\item as a catalyst: $\text{NO}_2$ speeds up the oxidation of atmospheric \textbf{sulfur dioxide} 二氧化硫 ($\text{SO}_2$) into $\text{SO}_3$, which then forms sulfuric acid in the rain.

\end{itemize}


\end{document}
